\documentclass[12pt, letterpaper, twoside]{article}
\usepackage[utf8]{inputenc}
\newcommand{\mycomment}[1]{}
\usepackage{imakeidx}
\usepackage{tcolorbox}
\usepackage{amsmath}

\title{RedHat Certified System Administrator}
\author{Jacopo Pela}
%\date{Luglio 2023}
\usepackage{hyperref}%
\makeindex
\begin{document}

\begin{titlepage}
\maketitle
\end{titlepage}

\tableofcontents
\newpage

\section{Linux Distributions}
Linux Distributions or distro are compleate system build around \textbf{linux kernel}\\
Linux distribution are differentieted by what sit arount the linux kernel\\
Different linux distro use different linux package manager to \textbf{install, update and remove} software form a centralized controlled source.
For RedHat we will talk about \textbf{rpm}.\\
Almost anything in a linux distribution is \textbf{Open Source}.\\
All linux distros share core linux tools like databases and programming languages.\\
RedHat follow the \textbf{File Hierarchy Standard (FHS)} for default file position.\\
RedHat releases new major versions every \textbf{three year}. Leap help you with a smooth transition from one version to the next one.
Version support come in three phases:
\begin{itemize}
	\item Full support for the first 5 years
	\item Maintenance support for the next 5 years
	\item Extended update support for selected security patches
\end{itemize}

\section{Introduction to the shell}
Commands are sunned on a shell and the default shell is bash.\\
Basic structure of a command is:\\

\centering{command <options> <arguments>}\\

Option begin with a dash ($-$) or a double dash ($--$)\\
Options influence how the command run, while argument are the things you want your command to work on\\

\section{Documentation}

\printindex
\end{document}
